% Transcription — fonds Alexandre Grothendieck, shelfmark 131, batch 2 (pages 21-40).
% Established from the facsimile archives/batches/131/batch-02.pdf (a mirror of
% https://grothendieck.umontpellier.fr/131.pdf), per the skill
% transcribe-grothendieck. The batch file carries no cover sheet: PDF page k
% is page 20+k of the folder (the Montpellier footer prints « 21 » ... « 40 »),
% and the archivists' pencil numbers agree wherever legible (« 22 », « 24 »,
% « 26 », « 28 », « 30 », « 32 », « 34 », « 36 », « 38 », « 40 »).
%
% Pass: Opus 5.5 (claude-opus-5-5), 2026-09-24 — first pass, unchecked against the pages by a human.
%
% Ten of the twenty pages are transcribed: the even pages 22-40, all in his
% hand. Absent: page 21, a blank leaf; pages 23, 25, 27, 29, 31, 33, 35, 37,
% 39, typed pages of an unrelated duplicated typescript (a Bourbaki draft,
% « n° 371 », on tensor products of algebras and the large algebra of a
% monoid) carrying nothing in his hand — someone else's text with no bearing
% on the folder's mathematics, skipped under the project's rule. Nothing in
% the batch is summarised. He did not paginate these leaves. The \section
% heading is ours, bracketed; « Proposition 1 », « Corollaire 1 »,
% « Corollaire 2 », « Proposition 2 », « Prop. 3 », « Lemme » are his.
%
% What the batch is. One run, undated: representability of a functor
% F : (Sch)°/S -> Ens by a prescheme locally of finite presentation (then
% locally quasi-finite) over S, from conditions checked on local rings.
%   22-24  Proposition 1: F of local nature (L_1), commuting with inductive
%          limits of rings (L_2), and (X, ξ) representing F on local
%          arguments essentially of finite type with finite residue extension
%          (R_loc) => (X, ξ) represents F; proof (injectivity, surjectivity
%          by spreading out u_t : Spec O_{T,t} -> X and gluing).
%   24-28  Corollaire 1 (S loc. noeth.): with L_3°° (F(B) -> F(B^) injective)
%          and L_4°° (F(B^) -> lim F(B_n) injective) and R_0 (representation
%          on artinian arguments), (X, ξ) represents F; proof through
%          T' = Spec B^ and the two projections of T' x_T T'. N.B.: the
%          natural statement is for a morphism G -> F. Corollaire 2: S
%          artinian, F representable by X loc. quasi-finite iff F is
%          prorepresentable by artinian local rings and satisfies L_1-L_4°°.
%   30-32  Proposition 2: A local noetherian complete, S' the punctured
%          spectrum, F_{S'} representable by X', formal part prorepresentable
%          by rings finite over A, and 4°) X'_1 -> X' an open immersion =>
%          F representable (glue X_1 and X'). Marginal programme: replace 4°)
%          by a checkable condition. Prop. 3: F_a representable by X_a over
%          every local S_a, the « modular » locus M(T, η) constructible,
%          L_1 and L_2 => F representable.
%   34-40  Definition of (T, η) « modulaire (loc.) » (a bijection onto the
%          η' whose residue value comes from k(t); fields of definition via
%          a diagram (Q)); N.B. on modular pairs being isomorphic; proof of
%          Prop. 3 by generisation + constructibility, a Lemme on
%          h_Z x_F h_Z giving an étale equivalence graph Z' ⇉ Z, trivial over
%          S_a, hence « presque modulaire » (X^(x), η^(x)), glued (via L_2).
%
% Notation fixed for the batch: his condition labels carry an arrow drawn
% under the L, rendered \underleftarrow{L} (\underrightarrow{L} where the
% arrow points right, L_4°° on pages 24 and 28); the two small circles on
% L_3, L_4 as ^{\circ\circ}; the completion as \hat{B}; the local ring as
% \mathcal{O}_{T,t}. Words he underlines are set in \emph. Marginal notes on
% pages 22, 30, 32, 38 are written slantwise and were read from rotated crops.
\documentclass[11pt,a4paper]{article}
% Le préambule commun à toutes les transcriptions du fonds Grothendieck.
%
% Il tient en une idée : **l'incertitude se compose, elle ne se résout pas.**
% Une transcription qui lisse un mot douteux a détruit la seule chose pour
% laquelle on la faisait — un lecteur doit pouvoir savoir, à chaque endroit,
% ce qui était sur le feuillet et ce qui a été ajouté. Les six macros qui
% suivent sont donc l'appareil critique, et non de la décoration.
%
% Les mêmes commandes sont reconnues par `scripts/render.mjs`, qui produit la
% vue de lecture du site. Une macro ajoutée ici doit l'être aussi là-bas,
% faute de quoi le rendu échouera bruyamment — ce qui est voulu : un rendu
% silencieusement faux est pire qu'un rendu absent.

\ProvidesPackage{grothendieck}[2026/08/08 Transcriptions du fonds Grothendieck]

\usepackage{amsmath,amssymb,amsthm}
\usepackage{mathrsfs}
\usepackage{stmaryrd}
% Les diagrammes commutatifs sont partout dans ce fonds : c'est la langue
% même de la géométrie algébrique de Grothendieck, et une transcription qui
% les rendrait en prose ne transcrirait rien.
\usepackage{tikz-cd}
% Français par défaut — c'est la langue des feuillets — mais l'anglais est
% chargé aussi : la traduction et le résumé sont des documents anglais, et la
% typographie française (espaces avant les deux-points) y serait une faute.
% Ces documents basculent par \selectlanguage{english} en tête de corps.
\usepackage[english,french]{babel}
\usepackage{fontspec}
\usepackage{geometry}
\geometry{a4paper,margin=25mm,marginparwidth=28mm}
\usepackage{marginnote}
\usepackage{xcolor}
% ulem, not soul. `soul`'s \st and \ul are built on a character-by-character
% scan that XeLaTeX's font machinery breaks: the document fails with an
% undefined control sequence rather than degrading. ulem does the same two jobs
% and is Unicode-engine safe. `normalem` stops it hijacking \emph, which the
% transcriptions use for Grothendieck's own emphasis.
\usepackage[normalem]{ulem}
\usepackage{hyperref}
\hypersetup{colorlinks=true,linkcolor=black,urlcolor=black,pdfborder={0 0 0}}

% --- Métadonnées du lot -----------------------------------------------------
% Reprises telles quelles par le script de rendu, qui les affiche en tête de
% la vue de lecture. Elles fixent ce que le fichier prétend couvrir : sans
% elles, une transcription flotte sans point d'attache dans un fonds de seize
% mille feuillets.

\newcommand{\folder}[1]{\gdef\grothFolder{#1}}
\newcommand{\batch}[1]{\gdef\grothBatch{#1}}
\newcommand{\leaves}[2]{\gdef\grothFirst{#1}\gdef\grothLast{#2}}
\newcommand{\foldertitle}[1]{\gdef\grothFoldertitle{#1}}
\gdef\grothFolder{?}\gdef\grothBatch{?}\gdef\grothFirst{?}\gdef\grothLast{?}\gdef\grothFoldertitle{}

% --- Le résumé --------------------------------------------------------------
%
% Il ouvre la lecture modernisée, sous le titre « Résumé », et il est composé
% en petit. Ce n'est pas un ornement : c'est le seul passage du document qui
% ne soit pas des mathématiques, et un lecteur doit pouvoir le reconnaître
% avant de l'avoir lu — puis le sauter s'il connaît déjà le sujet, ou s'y
% arrêter s'il ne le connaît pas. Le filet qui le suit marque la reprise du
% corps.

\newenvironment{resume}
  {\small\setlength{\parskip}{0.45em}}
  {\par\medskip\hrule\medskip}

% --- Mots-clés ---------------------------------------------------------------
%
% En anglais, en clôture du résumé de la lecture modernisée : le vocabulaire
% moderne sous lequel chercher ce que les feuillets construisent. Le manifeste
% du site (`npm run manifest`) les extrait pour étiqueter la cote entière —
% cette ligne est la source unique des tags, il n'y a pas de fichier à côté.
\newcommand{\keywords}[1]{\par\smallskip\noindent
  {\footnotesize\textsc{Keywords} --- \textit{#1}}\par}

% --- L'appareil critique ----------------------------------------------------

% Le numéro de feuillet, en marge. C'est l'ancre qui permet de régler un
% désaccord sur une formule en regardant *un* feuillet plutôt qu'une chemise
% de six cents. Le site s'en sert aussi pour tourner les pages du fac-similé
% quand on fait défiler la transcription.
\newcommand{\leaf}[1]{%
  \marginnote{\footnotesize\color{gray!70}\textsf{#1}}%
  \label{leaf:#1}\ignorespaces}

% Illisible : on ne devine pas. Un mot inventé qui se lit comme les autres est
% le pire résultat possible.
\newcommand{\ill}{\textcolor{red!60!black}{[\,\ldots\,]}}

% Lecture douteuse : le mot est donné, et signalé comme incertain.
\newcommand{\uncertain}[1]{\uline{#1}}

% Ajout de l'éditeur — une lettre manquante, un mot sous-entendu.
\newcommand{\add}[1]{\textcolor{blue!50!black}{[#1]}}

% Biffé par Grothendieck lui-même. Ce qu'il a rayé fait partie du document :
% une rature dit souvent où la pensée a changé de direction.
\newcommand{\struck}[1]{\sout{#1}}

% Note du transcripteur, sur l'état matériel du feuillet.
\newcommand{\note}[1]{\par\smallskip{\small\color{gray!60!black}\textit{[#1]}}\par\smallskip}

% Note marginale de Grothendieck, distincte du corps du texte.
\newcommand{\marginal}[1]{\par\smallskip\noindent\hspace{1em}%
  {\small\color{gray!60!black}#1}\par\smallskip}

% --- Titre ------------------------------------------------------------------

\AtBeginDocument{%
  \begin{center}
    {\large\bfseries Fonds Alexandre Grothendieck --- cote n\textsuperscript{o}~\grothFolder}\\[2pt]
    {\itshape\grothFoldertitle}\\[4pt]
    {\small feuillets \grothFirst--\grothLast\ (lot \grothBatch)}\\[2pt]
    {\footnotesize\color{gray!60!black}Transcription établie d'après le fac-similé numérisé
     par l'Université de Montpellier.\\
     Les lectures incertaines sont soulignées, les illisibles notées [\,\ldots\,],
     les ajouts entre crochets.}
  \end{center}
  \vspace{1em}}

\endinput


\folder{131}
\batch{2}
\pages{21}{40}
\dating{[vers 1971]}
\watermark{Édition de démonstration}
\foldertitle{Foncteurs représentables par schéma quasi fini : notes manuscrites (s.d.).}

\begin{document}

\section{[Critères de représentabilité par un préschéma localement de présentation finie]}
\note{titre de l'éditeur, entre crochets.}

\page{22}
\marginal{\uncertain{Conditions} \ill{} sur $F$ \ill{} ; \ill{} \ill{}, \uncertain{immersion} \ill{}.}
\note{note marginale à l'encre brune, écrite en biais dans le coin supérieur
gauche ; seuls quelques mots se lisent.}

\emph{Proposition 1} \struck{\uncertain{Lemme}}. $S$ préschéma,
$F : (\mathrm{Sch})^{\circ}_{/S} \to \mathrm{Ens}$, $X$ un préschéma loc. de
prés. finie sur $S$ \struck{$\in \operatorname{Ob} (\mathrm{Sch})_{/S}$},
$\xi \in F(X)$. On suppose
\note{« un préschéma loc. de prés. finie sur $S$ » est écrit au-dessus de la
ligne et renvoyé après $X$.}

$\underleftarrow{L}_1)$ $F$ de nature locale ;

$\underleftarrow{L}_2)$ $F$ « commute aux limites inductives d'anneaux » ;

$R_{\mathrm{loc}})$ $(X, \xi)$ représente $F$ quand on se restreint aux
arguments $T/S$ qui sont \struck{des anneaux} locaux et essentiellement de
type fini sur $S$, à extension résiduelle finie sur $S$.
\note{les trois étiquettes sont écrites sur des étiquettes antérieures,
biffées et illisibles ; le $L$ porte une flèche tracée dessous, rendue ici et
dans la suite par $\underleftarrow{L}$.}

Alors $(X, \xi)$ représente $F$, i.e. $\xi : X \to F$ est un isom.
\note{« $\xi : X \to F$ est un isom. » est au-dessus de la ligne.}

\emph{Dém.} Comme $F$ est de nature locale, on peut supposer $S$ affine
(\ill{}), alors $R_{\mathrm{loc}}$ signifie que pour tout spectre $T$ d'un
localisé $B$ d'une algèbre de t.f. sur \uncertain{un anneau} $A$,
\[
  \operatorname{Hom}_S(\struck{\operatorname{Spec}(B)}\, T, X)
  \xrightarrow{\ \sim\ } F(T)
\]
par $\xi$. \struck{\ill{}} On est \uncertain{ramené} à : prouver que ceci
reste vrai pour tout $T$ \struck{loc.} de t.f. sur $S$ (\ill{}).
\note{« spectre $T$ d'un » est au-dessus de la ligne ; le même mot illisible
entre parenthèses suit « $S$ affine » et « sur $S$ ».}

1°) \emph{Injectivité}. Soient $u, v : T \rightrightarrows X$
\struck{définis} tels que $F(u)(\xi) = F(v)(\xi)$, prouvons $u = v$. Par
l'hyp. $R_{\mathrm{loc}}$, $u$ et $v$ ont même restriction aux
$\operatorname{Spec} \mathcal{O}_{T,t}$ (à ext. résiduelle finie), d'où la
conclusion.

2°) \emph{Surjectivité}. Soit $\eta \in F(T)$. Pour tout $t \in T$ (à ext.
résiduelle finie), il existe grâce à $R_{\mathrm{loc}}$ un
$u_t : \operatorname{Spec} \mathcal{O}_{T,t} \to X$
(\uncertain{unique par 1°)}) induisant
$\uncertain{\eta_t} \in F(\operatorname{Spec} \mathcal{O}_{T,t})$.
\struck{En} outre
\note{« (à ext. résiduelle finie) » est ajouté dans la marge droite et relié
par une accolade aux deux lignes. La lettre de $\eta_t$ est tracée comme un
$\xi$.}

\page{24}
\struck{en vertu de b)}, comme $X$ est loc. de prés. finie sur $S$, $u_t$
provient d'un \struck{$A \to \varinjlim B_f$}
$v_t : U_t \to X$, où $U_t$ est un voisinage ouvert de $t$. En vertu de b),
prenant $U_t$ assez petit, on aura $F(v_t)(\xi) = \eta | U_t$.
En vertu de 1°), on a, pour $t, t' \in T$,
\[
  v_t | U_t \cap U_{t'} = v_{t'} | U_t \cap U_{t'} ,
\]
où les $U_t$ recouvrent $T$, \ill{} un $v : T \to X$ \struck{induisant}
induisant les $v_t$.
Alors en vertu de a), on a $F(v)(\xi) = \eta$.

CQFD
\note{les renvois « a) » et « b) » visent vraisemblablement les conditions de
la page 22 sous leurs premières étiquettes, biffées. « où les $U_t$
recouvrent $T$ » est au-dessus de la ligne.}

\marginal{\struck{Corollaire 1}}
\emph{Corollaire 1}. Sous les conditions préliminaires de la
\uncertain{prop.}, supposons en plus $S$ loc. noeth., \uncertain{conditions
ci-dessus} \struck{Supposons aussi les conditions a), b), c) de} conditions
suivantes \struck{de} : pour tout anneau local $B$ ess. de type fini sur $S$,

$\underleftarrow{L}_3^{\circ\circ})$ \struck{a) $F$ \ill{} compatible \ill{}
\ill{}} $F(B) \to F(B')$ \uncertain{injectif} \struck{\ill{} complétion
$B'$, \ill{} et \ill{} complets} [il \struck{existe} \ill{} dans certains
cas \ill{}]
\marginal{$F(B) \to F(\hat{B}) \to \varprojlim$ \ill{}, \emph{injectif}}
\note{passage très raturé : la condition, d'abord posée sur un complété $B'$,
est biffée ligne par ligne ; la note de marge, reliée par une flèche, en
donne la forme retenue.}

$\underrightarrow{L}_4^{\circ\circ})$ \struck{a) $F$ \ill{}} pour tout anneau
local $B$ (à ext. résiduelle finie sur $S$) ess. de type fini sur $S$,
\[
  F(\hat{B}) \xrightarrow{\ \mathrm{can}\ } \varprojlim F(B_n)
\]
est \emph{injectif}.

$R_0)$ $(X, \xi)$ représente $F$ pour des arguments artiniens ess. de type
fini sur $S$, à ext. résiduelle finie sur $S$.

Si $\underleftarrow{L}_1$, $\underleftarrow{L}_2$,
$\underleftarrow{L}_3^{\circ\circ}$, $\underrightarrow{L}_4^{\circ\circ}$
\struck{\ill{}} $R_0$ sont satisfaits, alors $F$ est représentable par
$(X, \xi)$.
\note{l'énoncé du corollaire est tenu par un trait vertical dans la marge
gauche.}

\page{26}
Il suffit de prouver que $F$ satisfait à $R_{\mathrm{loc}}$ \struck{\ill{}}.
On peut supposer $S$ affine, soit $T$ loc. ess. de t. fini sur $S$, défini
par $B$, $T_n = \operatorname{Spec}(B_n)$, $T' = \operatorname{Spec} \hat{B}$,
$\eta' =$ image de $\eta$ par $F(T) \to F(T')$. À prouver que
\uncertain{$(*)$} est un isom.
\note{« on peut supposer $S$ affine » est au-dessus de la ligne ; « image de
$\eta$ par $F(T) \to F(T')$ » est ajouté dans la marge droite. Le signe
cerclé renvoie à une application qui n'est pas écrite sur ce feuillet,
vraisemblablement $\operatorname{Hom}_S(T, X) \to F(T)$.}

1°) C'est injectif. En effet, il suffit de \uncertain{noter} que si
$u, v : T \to X$ sont tels que $F(u)(\xi) = F(v)(\xi)$, on a
$F(u_n)(\xi) = F(v_n)(\xi)$ [induites sur $T_n$ par $u$, $v$], d'où par
$R_0$ $u_n = v_n$ pour tout $n$, d'où $u = v$ car $B$ noeth. donc
$B \subset \hat{B}$.

2°) C'est surjectif. \struck{Il suffit} Soit $\eta \in F(T)$. Il existe par
$R_0$ des $u_n : T_n \to X$ qui induisent les $\eta_n$, uniques, donc
compatibles, d'où \struck{\ill{}} $u' : T' \to X$ \struck{unique} induisant
les $u_n$, d'où \struck{\ill{}} en vertu de $\underleftarrow{L}_{\mathrm{loc}}$
\struck{la relation} $F(u')(\xi) = \eta'$.

Je dis qu'on a \struck{\ill{}} $u' p_1 = u' p_2$. En vertu de 1°), cela
\struck{\ill{}} signifie
$F(u' p_1)(\xi) = F(u' p_2)(\xi)$, qui se réduit à
$F(p_1)(\eta') = F(p_2)(\eta')$, qui résulte de $\eta' = F(f)(\eta)$ et
$f p_1 = f p_2$.
\[
  T' \times_T T' \;\overset{p_1}{\underset{p_2}{\rightrightarrows}}\; T'
  \xrightarrow{\ f\ } T ,
  \qquad u' : T' \to X
\]
\note{diagramme dans la marge gauche ; $u'$ y est une flèche oblique de $T'$
vers $X$.}

\page{28}
Donc il existe $u : T \to X$ tel que $u' = u f$. Alors on a \struck{\ill{}}
$\eta = F(u)(\xi)$, car en vertu de $\underleftarrow{L}_3$ a) il suffit de
vérifier que $F(f)(\eta) = F(f) F(u)(\xi)$, i.e. $\eta' = F(u')(\xi)$, ce
qui est vrai.
\note{les deux $F(f)$ portent un trait qui peut être un accent : $F(f')$.}

OK

\emph{N.B.} Il est plus naturel d'énoncer ces deux résultats ainsi :
\ill{} \ill{} un \uncertain{hom.} de foncteurs $G \to F$ satisfaisant les
conditions envisagées pour $F$, et \ill{} \ill{} \uncertain{qu'il est} un
\emph{isom.}

\emph{Corollaire 2}. Soient $S$ schéma artinien,
$F : (\mathrm{Sch})^{\circ}_{/S} \to (\mathrm{Ens})$, pour que $F$ soit
représentable par un $X$ loc. quasi fini sur $S$, il faut et il suffit que
$F$ satisfasse les conditions suivantes :

1°) $F$ est proreprésentable, les anneaux locaux qui interviennent dans la
proreprésentation de $F$ sont \struck{\ill{}} \emph{artiniens}, i.e. finis
sur $S$.

2°) $F$ satisfait les conditions $\underleftarrow{L}_1$,
$\underleftarrow{L}_2$, $\underleftarrow{L}_3^{\circ\circ}$,
$\underrightarrow{L}_4^{\circ\circ}$.

\page{30}
En effet, grâce à 1°) et à $\underleftarrow{L}_1$ (\ill{}), on dispose d'un
$(X, \xi)$. On peut lui appliquer le cor. 1, d'où la conclusion.

\struck{Prop.} Plus généralement, on obtient

\struck{Corollaire 3}

\emph{Proposition 2}. $S = \operatorname{Spec}(A)$, $A$ anneau local noeth.
\emph{complet}, $S'$ = complémentaire de l'origine, \struck{\ill{}}
$F : (\mathrm{Sch})^{\circ}_{/S} \to (\mathrm{Ens})$. On suppose que

1°) $F_{S'}$ représentable par un \struck{sch.} préschéma $X'$ loc. de type
fini sur $S$.

2°) La restriction de $F$ aux $B$ artiniens finis sur $S$ est
proreprésentable, les anneaux locaux qui interviennent dans la
proreprésentation \struck{étant finis sur} sont \emph{finis sur $A$}.
\note{au-dessus de la ligne : « il suffit de le vérifier pour $A \to$
\ill{} », la fin raturée.}

3°) $F$ satisfait les conditions $\underleftarrow{L}_1$
$\underleftarrow{L}_2$ $\underleftarrow{L}_3^{\circ\circ}$
$\underleftarrow{L}_4^{\circ\circ}$.
\note{ici la flèche sous le quatrième $L$ paraît tournée vers la gauche.}

On construit grâce à 2°) un
$X_1^{*} = \coprod_i \operatorname{Spec}(B_i)$ et
$\xi_1^{*} \in F(X_1^{*})$. Considérons $X'_1 = X_1 | S'$,
$\xi'_1 = \xi_1 | X'_1$, d'où $\xi'_1 = F(u)(\xi')$ où
$u : X'_1 \to X'$ est un morphisme. On suppose

4°) $u$ est une \emph{immersion ouverte}.

\marginal{\ill{} : supposons donnés $X_U$ sur $U$ et $Z$ sur $S$,
$X_U \xrightarrow{\sim} F_U$, $Z \to F$, tels que (i) $Z_U \to X_U$ une
immersion ouverte ; (ii) \struck{\ill{}} $Z \to F$ \ill{} en les pts de
$Z_0$ ; (iii) $Z_0 \to F_0$ un isom.}
\note{note écrite en biais dans le coin inférieur gauche, d'une autre
plume ; elle reformule les hypothèses de la proposition.}

\page{32}
Sous ces conditions, $F$ est représentable.

En effet, on peut recoller $X_1$ et $X'$ suivant $u : X'_1 \to X'$, on
obtient $X$, et sur $X$, grâce à $\underleftarrow{L}_1$, on a un
$\xi \in F(X)$, induisant $\xi_1$ et $\xi'$. On lui applique le
corollaire 1, qui implique que $(X, \xi)$ représente $F$.
\marginal{Transformer 4°) en une \ill{} plus facile à vérifier, avec une
bonne \ill{} de \uncertain{dim} 1.}

\emph{Prop. 3}. Soit $F : (\mathrm{Sch})^{\circ}_{/S} \to \mathrm{Ens}$, on
suppose : \struck{\ill{}}

1°) Pour tout $a \in S$, posant
$S_a = \operatorname{Spec} \mathcal{O}_{S,a}$, le foncteur $F_a = F_{/S_a}$
est représentable par un préschéma $X_a$ loc. de t.f. sur $S_a$.

2°) Pour tout $T$ sur $S$, loc. de type fini sur $S$, et $\eta \in F(T)$,
l'ensemble $M(T, \eta)$ des $t \in T$ tels que
$\eta_t \in F(\uncertain{T_t})$ soit \emph{modulaire} (loc.), est une partie
\emph{constructible} de $T$.

3°) $F$ satisfait $\underleftarrow{L}_1$ et $\underleftarrow{L}_2$.

Sous ces conditions, $F$ est représentable.
\marginal{Transformer en une condition de modularité « générique ».}
\note{3°) est écrit en bas du feuillet et renvoyé par une flèche avant la
conclusion ; « $M(T, \eta)$ » et « (loc.) » sont ajoutés au-dessus de la
ligne. La note de marge est reliée par une flèche à 2°).}

\page{34}
\emph{N.B.} On dit que, pour un $T$ \emph{local} sur $S$ et
$\eta \in F(T)$, $(T, \eta)$ est \emph{modulaire} (loc.) si :

1°) pour tout $T'$ local sur $S$, l'application
\[
  \operatorname{Hom}_{\mathrm{loc}, S}(T', T) \longrightarrow F(T')
\]
définie par $\eta$ est une \struck{application} bijection du premier
ensemble sur le sous-ensemble des $\eta' \in F(T')$ tels que
$\eta'(t') \in F(k(t'))$ soit dans
\[
  \operatorname{Im}\bigl[\operatorname{Hom}_k(k(t'), k(t))
  \longrightarrow F(k(t'))\bigr] ,
\]
où la deuxième flèche est définie par $\eta(t) \in F(k(t))$.
\marginal{$t$, $t'$ les pts fermés de $T$, $T'$}

2°) Pour tout diagramme commutatif de corps \uncertain{(Q)}, où
$L = k(t)$, $L' = k(t')$, $k \to L$ et $k \to L'$ les homomorphismes
canoniques, et tout $\beta' \in F(L')$ tel que
$F(j')(\beta') = F(j)(\beta)$ [où $\beta = \eta(t)$], on a
\[
  \begin{cases}
    j'(L') \supset L & (\text{d'où } \lambda : L \to L') \\
    \beta' = F(\lambda)(\beta)
  \end{cases}
\]

\begin{tikzcd}
L \arrow[r, "j"] \arrow[dr, dashed, "\lambda"] & M \\
k \arrow[u, "i"] \arrow[r, "i'"'] & L' \arrow[u, "j'"']
\end{tikzcd}

\note{le diagramme (Q) est en marge gauche, avec l'annotation
$\eta(t) = \beta$ devant $L$ et $\beta' = $ \ill{} sous $L'$ ; la flèche
$\lambda$ est pointillée.}

[\emph{N.B.} On \struck{dit que \ill{} que} exprime 2°) en disant que $L$ est
un \emph{corps de définition} pour $\beta$.]

\begin{tikzcd}
B \arrow[r] & L & \\
 & k \arrow[u] \arrow[r] & L' \\
A \arrow[uu] \arrow[ur] \arrow[rr] & & B' \arrow[u]
\end{tikzcd}

\note{croquis en marge gauche, sous les indications $t \in T$, $s \in S$,
$T' \ni t'$ séparées par une courbe ; plusieurs lettres y sont biffées
(un $k$ à gauche de $B$ et de $A$, un objet à droite de $B'$), et une flèche
pointillée descend vers $L'$ depuis une lettre biffée voisine de $L$ ; on n'a
retenu que les flèches non biffées.}

\page{36}
Si $a$ est le pt de $S$ au-dessous de $t$, et si $F_a$ est
\emph{représentable} par $X_a$ \struck{\ill{}} loc. de type fini sur $S_a$,
\struck{\ill{} fini \ill{}} \ill{} $\eta$ définit $u : T \to X_a$, posant
$x = u(t)$, d'où $\bar{\eta} : \mathcal{O}_{X_a, x} \to \mathcal{O}_{T,t}$,
le fait que $\eta$ soit modulaire (loc.) signifie que $\bar{\eta}$ est un
\emph{isomorphisme}.
\note{la lettre surmontée d'une barre est tracée comme $\eta$.}

\struck{Rem.} \emph{N.B.} Si $\eta \in F(T)$, $\eta' \in F(T')$ sont
modulaires (loc.), et s'il existe un diagramme commutatif (Q) tel que
$F(j)(\eta(t)) = F(j')(\eta'(t'))$, alors (\uncertain{et alors seulement})
il existe un \uncertain{(unique)} $\varphi : T \to T'$ \struck{tel} que
$\eta = F(\varphi)(\eta')$.

\emph{Preuve prop. 3} (on peut supposer $S$ noeth.). De 1°) résulte que
l'ensemble \struck{\ill{}} $M(T, \eta)$ est stable par \struck{\ill{}}
générisation. Par 2°), il est donc \uncertain{ouvert}. \struck{Donc}
Utilisant $\underleftarrow{L}_2$, (pour tout $a \in S$ et $x \in X_a$ sur
$a$,) \struck{\ill{}} il existe \ill{} un \ill{} $X^{(x)}$ \struck{\ill{}}
\uncertain{noethérien}, \struck{\ill{}} de t.f. sur $S$, et
$\eta^{(x)} \in F(X^{(x)})$, modulaire (loc.) en
\note{« on peut supposer $S$ noeth. » est ajouté au-dessus de la ligne ;
« Utilisant $\underleftarrow{L}_2$ » est interlinéaire.}

\page{38}
tous les pts de $X^{(x)}$.
\struck{Pour tout \ill{} \ill{} \ill{} que les $X^{(x)}$ se recollent \ill{}.
Pour tout couple $(x, y)$ correspondant à $(a, b)$, soit $X^{(a,b)}$
l'ensemble des $u \in X^{(a)}$ qui sont apparentés à un pt de $X^{(b)}$.
D'après \ill{} $\underleftarrow{L}_2$, \ill{} résulte \ill{} \ill{}
\ill{} \ill{}. \ill{} De plus, on a un}
\note{tout ce passage est encadré et barré de longs traits obliques.}

\emph{Lemme}. Soit $(Z, \eta)$ ($Z$ loc. de t.f. sur $S$,
$\eta \in F(\uncertain{S})$) tel que pour tout $z \in Z$, $\eta_z$ soit
modulaire. Considérons $h_Z \xrightarrow{\eta} F$, \struck{et} alors
$h_Z \times_F h_Z \subset h_{Z \times Z}$ est représentable par un
sous-préschéma de $Z \times Z$, soit $Z'$. Les projections
$Z' \overset{\mathrm{pr}_1}{\underset{\mathrm{pr}_2}{\rightrightarrows}} Z$
sont des isomorphismes locaux, [et font de $Z'$ un \struck{préschéma} graphe
d'équivalence dans $Z$ \ill{}].
\note{on attendrait $\eta \in F(Z)$.}
\marginal{Reprendre les conditions \uncertain{(nécessaires)} sur $F$.
Celles \ill{} : \ill{} puis \ill{} …}

AQT
\note{abréviation non résolue, centrée sous le lemme.}

Or sous les conditions préliminaires, prenant $Z = X^{(x)}$, on voit que
quand on se restreint au-dessus de $S_a$,
$Z'_a \rightrightarrows Z_a$ est la relation d'équivalence \emph{triviale}.
Donc

\page{40}
il en est de même de $Z' | U \rightrightarrows Z | U$, où $U$ est un
voisinage ouvert convenable de $a$. En d'autres termes, \emph{à condition de
prendre $X^{(x)}$ assez petit}, on peut supposer que
$(X^{(x)}, \eta^{(x)})$ est « \uncertain{presque} \emph{modulaire} »
[Précision de la signification de ce terme :
$\eta_x : h_{X^{(x)}} \to F$ est une immersion ouverte]. Alors, une autre
proposition « formelle » nous permet de recoller les $X^{(x)}$ et de trouver
un $(X, \xi)$ qui représente $F$ (grâce à $\underleftarrow{L}_2$).
\note{« presque » est cerclé ; la précision entre crochets est ajoutée
au-dessus de la ligne et reliée par une accolade.}

\end{document}
